\scanpage{187}
\mspage{136}
la seule différence avec celui formulé à la p.~\deleted{\unsure{126|120}}\unsure{128|126} étant la direction opposée des flèches $\lambda_1,\lambda_2$. Pour \unsure{tt} carré commutatif
\diagram{\begin{tikzcd}[ampersand replacement=\&,row sep=small,column sep=small]
\&\&X_0\arrow[dl,"\alpha_1"']\arrow[dr,"\alpha_2"]\&\\
(\ast)\&X_1\arrow[dr,"f_1"']\&\&X_2\arrow[dl,-,"f_2"]\\
\&\&Y\&{,}
\end{tikzcd}}
on trouve encore une flèche canonique
\[
\widetilde{\unsure{\varphi|\rho}}:\widetilde{\mathscr X}\longrightarrow Y .
\]
Il serait \deleted{\unsure{tout}} aussi raisonnable de dire que le carré $(\ast)$ est $W$-homotopiquement cocartésien, lorsque $\widetilde\varphi$ est $\in W$, que d'adopter la définition 3, page 131 -- on trouverait exactement les $=\!{}^{\mathrm{s}}$ corollaires. Les deux définitions seraient-elles équivalentes ??

Je note que $\widetilde{\mathscr X}$ étant cofibrée sur $\Psi=\Phi^0$, $\widetilde{\mathscr X}^0$ est fibrée sur $\Psi^0=\Phi$. De façon précise, on aura
\diagram{\begin{tikzcd}[ampersand replacement=\&,row sep=tiny,column sep=small]
\&X_1^0\\
\widetilde{\mathscr X}^{0}=\int (X_{\unsure{0}})^0\arrow[ru,"\alpha_1^0"]\arrow[rd,"\alpha_2^0"']\&\\
\&X_2^0\rlap{,}
\end{tikzcd}}
% archive sheet number: 187
